\section{Results}
\label{sec:results}

We report results of evaluating the generated one-shot
\texttt{gpt-4o-mini-2024-07-18} unit-test suites on 120 functions (60 Java from 8
Defects4J subject programs; 60 Python from \texttt{flask}/\texttt{requests}). Coverage
figures are computed with JaCoCo (Java) / \texttt{coverage.py} (Python); mutation score
with PIT (Java) and, for Python, \texttt{mutmut} for the LLM suites and a custom AST-level
engine for the Pynguin baseline (Section~\ref{sec:method}), all at the $\alpha = 0.05$
significance level. For matched-pairs analysis we report the rank-biserial $r$ (Wilcoxon
for matched pairs); for the coverage--complexity comparison (RQ3) we report Spearman
$\rho$. Per the protocol, any suite that failed to compile or run is scored 0\% coverage
and marked \textsc{invalid}. We report metrics at two levels side by side: \emph{all}
($N=60$ per language, invalid $=0$) and \emph{effective} ($>0\%$ coverage). The effective
figures should be read as conditioned on validity and only in the context of the
all-level results.

\subsection{Test validity (a primary finding)}
Across Python, only 32/60 (53.3\%) of the generated suites executed even one test; the
other 28/60 (46.7\%) were \textsc{invalid}. Of the 32 valid Python suites, only 23/60
(38.3\%) exercised the target function (had $>0\%$ coverage); the other 9/60 (15.0\%) were
mocks or self-implementations, or failed to launch. Breaking down the 28 invalid Python
suites:
\begin{itemize}
  \item \textbf{23 --- \texttt{ImportError}} (LLM-attributable): the test tries to import
        a symbol that is not present in the target module (e.g.\ an instance method
        imported as a module-level function). A few appear data/platform-specific
        (\texttt{proxy\_bypass\_registry} is Windows-only; \texttt{merge\_environment\_settings}
        is a method, not module-level).
  \item \textbf{5 --- \texttt{TypeError}/\texttt{AttributeError}} (environment-dependent):
        possibly tied to \texttt{flask}/\texttt{requests} versions, to be confirmed in the
        canonical environment.
\end{itemize}
Python's 46.7\% invalidity is the first key finding, and the distribution shows the
problem is not mainly environment-driven: 23 of the 28 failures are LLM-attributable, and
only 5 are environmental. The Python failures are not random --- \texttt{ImportError}
dominates, driven by importing an instance method as if it were a module-level function.

Java's invalidity looks rarer on the surface: only 9/60 (15\%) of Java suites were scored
\textsc{invalid} by our pipeline. That figure understates the problem, and a post-hoc audit
(Section~\ref{sec:threats}) shows why: the \texttt{compiled} flag was inferred from a
function's presence in the JaCoCo report rather than measured, and JaCoCo lists every class
in a project whether or not it was ever loaded. Compiling each suite directly with
\texttt{javac} against the project classpath gives 17/60; the remaining 43 fail to compile,
most often on ambiguous overloads or on targets that are \texttt{private}. Either way the
downstream conclusion is unchanged, because it rests on the effective count: only 5 suites
actually exercised the target method. Java's real problem was wrong-target invocation: the LLM wrote a call to a method
that was not the target but was present and accessible. For instance,
\texttt{CommandLine.getOptionValues(\ldots)} was called when the sampled function was
actually the nested \texttt{CommandLine.Builder.getOptionValues(\ldots)}; elsewhere a
same-named method inherited from a parent class was invoked instead of the intended
override. Java's suites were syntactically and semantically valid but did not target the
intended code --- the dominant failure mode, and a severe one, since they provided no
coverage of the target. After validity, only 5/60 (8.3\%) of Java suites were effective,
compared to 23/60 (38.3\%) for Python. Of the 9 suites our pipeline scored \textsc{invalid},
6 targeted \texttt{AccurateMath} --- a class packed with overloaded methods of similar
signatures --- and 3 targeted \texttt{Gson}'s generic adapters (\texttt{getAdapter},
\texttt{getDelegateAdapter}, \texttt{fromJson}); only one of the 9 targeted a
\texttt{private} method. Java's validity problems cluster in classes with overloaded
methods and generics, not around private methods --- which shapes the fix: making the
target \texttt{public} would help in at most 1 of the 9 cases, whereas a class-level
view (the methods available across the class, its hierarchy, and their names) would be far
more informative. Section~\ref{sec:rq4} tests exactly this hypothesis.

\subsection{RQ1 --- Branch coverage vs.\ 80\%}
\emph{RQ1: When \texttt{gpt-4o-mini} is given a single one-shot prompt and asked to write a
unit test for one individual function of medium cyclomatic complexity ($5 \le CC \le 10$),
does the suite it produces achieve a median branch coverage of at least 80\,\% over that
function's own line range?}

Table~\ref{tab:rq} shows the coverage figures. Across all 60 Python functions, the median
coverage was \textbf{0\%} (IQR 0.0--70.0): for at least half the sample, the suite covered
none of the target function's branches. Only 23 of the 60 suites exercised any target code,
and their effective coverage had a median of \textbf{75.0\%} ($n=23$, IQR 64.6--95.8; 11 of
the 23 reached $\ge 80\%$, and 6 reached 100\%). That effective median is still below the
bar, and the one-sample Wilcoxon test did not support the threshold ($p=0.67$, $r=-0.10$,
negligible). Java followed the same pattern one level lower: all-function median
\textbf{0\%} (IQR 0.0--0.0), effective median \textbf{50.0\%} on a much smaller subset
($n=5$, IQR 50.0--87.5, the 8.3\% effective rate; $p=0.91$, $r=-0.60$, large). Pooled across
languages, nothing changed: $n=120$, all-function median \textbf{0\%}; pooled effective
median \textbf{75.0\%} ($n=28$; $p=0.84$, $r=-0.22$, small). Therefore, we fail to reject
$H_0$: \textbf{H1 (median branch coverage $\ge 80\,\%$) is not supported} at either level, in
either language, and every effective median sat below 80\% (Fig.~\ref{fig:dist}). Those
effective figures are best-case: they are conditioned on the model managing to compile and
invoke the intended method. The most honest figure is the all-level 0\% median.

\subsection{RQ2 --- Mutation score}
\emph{RQ2-A: Does the same suite detect injected faults at a median mutation score of at
least 60\,\%? RQ2-B: Matched function by function, does it achieve a higher mutation score
than the established automated generators?}

\textbf{(A) vs.\ 60\%.} The pattern repeats: across all 60 functions the median score was
\textbf{0\%} (IQR 0.0--0.0), with only 17/60 functions producing a suite with a measurable
mutation score (Table~\ref{tab:rq}). Restricting to those 17 green suites, the effective
median rose to \textbf{36.84\%} ($n=17$, IQR 13.6--55.0; only 4/17 cleared 60\%), and the
one-sample Wilcoxon did not support the threshold ($p=0.99$, $r=-0.66$, large). Java was
comparable: all-function median \textbf{0\%} ($n=60$, IQR 0.0--0.0), and even restricting to
the 54 functions where PIT produced a score, the effective median was \textbf{0\%}
(one-sample Wilcoxon, $p=1.00$, $r=-1.00$, large). Nearly all Java-generated suites failed
the most basic mutation analysis, killing almost no mutants (Fig.~\ref{fig:dist}). Pooled
across the 71 functions with a defined mutation score, the median was likewise \textbf{0\%}
($p=1.00$, $r=-0.97$, large). Therefore, we fail to reject $H_0$ in either language:
\textbf{H2-A (median mutation score $\ge 60\,\%$) is not supported} --- the generated tests
were ineffective at killing mutants.

\textbf{(B) vs.\ baseline.} Paired Wilcoxon tests --- gpt-4o-mini against each baseline,
matched by \texttt{function\_id} on mutation score --- show gpt-4o-mini substantially
worse than the two automated Java baselines: EvoSuite ($n=54$ pairs, 16 non-tied;
$p=0.00044$, $r=-1.00$, large) and Randoop ($n=54$ pairs, 11 non-tied; $p=0.0033$,
$r=-1.00$, large). Therefore \textbf{H2-B is rejected for Java}: the baselines score above
the model, not below it. Against the Python baseline Pynguin the direction
flips: gpt-4o-mini is considerably better ($n=13$, $p=0.010$, $r=+0.85$). In other words,
for Java both EvoSuite and Randoop killed significantly more mutants than gpt-4o-mini's
tests, while for Python gpt-4o-mini killed more than Pynguin (Fig.~\ref{fig:rq2b}). That
Python win should be read with caution: Pynguin was severely underpowered on this dataset
(Section~\ref{sec:threats}) and the tests it produced barely killed any mutants ---
Pynguin's mutation score was below 60\% for the vast majority of functions, and its median
across the 13 compared functions was 0\%. So gpt-4o-mini beat Pynguin largely because
Pynguin sat so far below the threshold. \textbf{H2-B therefore holds only against Pynguin},
and for a reason that says more about the baseline than about the model.

\begin{figure}[t]
\centering
\includegraphics[width=\linewidth]{figures/fig2_comparison.png}
\caption{Mutation score, gpt-4o-mini vs.\ EvoSuite, matched by function ($n=54$ paired
Java functions). EvoSuite wins in every pair where the two differ ($r=-1.00$, $p=0.00044$).}
\label{fig:rq2b}
\end{figure}

\subsection{RQ3 --- Complexity vs.\ quality}
\emph{RQ3: Within the narrow 5--10 band alone, does cyclomatic complexity correlate
negatively with the branch coverage the model achieves?}

For the Python effective subset, the Spearman rank-order correlation was
$\rho = +0.18$ (negligible), not significant at $\alpha=0.05$ ($p=0.41$, $n=23$): there was
no negative relationship between complexity and coverage inside the CC 5--10 window. Java's
effective subset ($n=5$) gave an even weaker $\rho = +0.13$ ($p=0.83$), and across subsets
($n=28$) the correlation was $\rho = +0.16$ ($p=0.40$). For neither language could $H_0$
be rejected at the 0.05 level, so complexity within the CC 5--10 window was not a
significant predictor of coverage quality. The pre-registered decision rule for this
hypothesis was $\rho < -0.5$, which none of the observed values satisfied. Therefore, we fail
to reject $H_0$: \textbf{H3 (Spearman $\rho < 0$) is not supported} --- if anything, the
results suggest a positive (but non-significant) relationship between CC and coverage.

\begin{table*}[t]
\centering
\caption{Per-RQ summary (gpt-4o-mini). ``all'' $=60$/lang.\ (invalid/no-touch $=0$); ``eff.''
$=$ effective subset (compiled \& touches target); ``meas.'' $=$ subset with a measurable
mutation score.}
\label{tab:rq}
\begin{tabular}{lllrrl}
\hline
RQ & Metric & Level & Median & $p$ & Effect \\
\hline
RQ1 Py & Coverage$\ge$80 & all (60) & 0\% & -- & -- \\
       &                 & eff (23) & 75.0\% & 0.67 & $r{=}{-}0.10$ \\
RQ1 Java & Coverage$\ge$80 & all (60) & 0\% & -- & -- \\
         &                 & eff (5) & 50.0\% & 0.91 & $r{=}{-}0.60$ \\
RQ1 All & Coverage$\ge$80 & all (120) & 0\% & -- & -- \\
        &                & eff (28) & 75.0\% & 0.84 & $r{=}{-}0.22$ \\
RQ2A Py & Mutation$\ge$60 & all (60) & 0\% & -- & -- \\
        &                & meas.\ (17) & 36.84\% & 0.99 & $r{=}{-}0.66$ \\
RQ2A Java & Mutation$\ge$60 & all (60) & 0\% & -- & -- \\
          &                 & meas.\ (54) & 0\% & 1.00 & $r{=}{-}1.00$ \\
RQ2B Java & GPT vs EvoSuite & paired (54) & -- & $<$0.001 & $r{=}{-}1.00$ \\
RQ2B Java & GPT vs Randoop & paired (54) & -- & 0.003 & $r{=}{-}1.00$ \\
RQ2B Py & GPT vs Pynguin & paired (13) & $\Delta$med $+$36.84pp & 0.010 & $r{=}{+}0.85$ \\
RQ3 Py & CC$\sim$Cov & eff (23) & $\rho{=}{+}0.18$ & 0.41 & $+0.18$ \\
RQ3 Java & CC$\sim$Cov & eff (5) & $\rho{=}{+}0.13$ & 0.83 & $+0.13$ \\
RQ3 All & CC$\sim$Cov & eff (28) & $\rho{=}{+}0.16$ & 0.40 & $+0.16$ \\
\hline
\end{tabular}
\end{table*}

\begin{figure}[t]
\centering
\includegraphics[width=\linewidth]{figures/fig1_distribution.png}
\caption{Branch coverage (left) and mutation score (right) of gpt-4o-mini across all 60
functions per language ($N=120$); the median is 0\% in every panel (dashed: 80\% / 60\%
thresholds). Effective-subset medians (75.0\% Python / 50.0\% Java coverage; 36.84\% Python
mutation) are given in the text.}
\label{fig:dist}
\end{figure}

\subsection{Summary of findings}
Neither the 80\% coverage nor the 60\% mutation threshold is met in either language
($N=120$). Even in the most favourable group, the effective subset, coverage is 75.0\%,
still below 80\%. When counting all 120 functions, the median is 0\%. There is no clear
link between complexity and quality within CC 5--10 for either language. What matters
most here is test validity, and it breaks differently in each language: Python mostly
dies at compile time (46.7\% \textsc{invalid}, dominated by \texttt{ImportError}), while
Java mostly passes our pipeline's compile check and then measures nothing (46/60 $=$ 76.7\%
were scored \texttt{compiled} yet touch 0\% of the target; the direct \texttt{javac} audit
of Section~\ref{sec:threats} puts genuine compilation lower still), through silent
wrong-target/overload resolution. Effective suites make up only 8.3\% of the Java
sample and 38.3\% of the Python one. On paired comparisons, gpt-4o-mini loses decisively to
both automated Java baselines (EvoSuite, Randoop; $r=-1.00$ each) but beats the Python
baseline Pynguin, whose own coverage on this dataset sits at 0.0\% because of
tool--environment issues rather than any real task difficulty (Table~\ref{tab:rq}).
